\abstract{%
\linespacing{1.2}
Machine learning models in production can lose accuracy without warning when the distribution of their input data drifts. Unsupervised drift detectors can signal that drift has occurred, but D3, a discriminative detector recommended by a recent benchmark, relies on an uncalibrated AUC threshold, compares two short adjacent windows, and gives no indication of which features changed. This thesis proposes Enhanced Adversarial Drift Detection (EADD), an integrated framework that combines a LightGBM classifier two-sample test, a reservoir-sampled reference window, prediction entropy monitoring and an adaptive sequential permutation test (ASPT), and that ranks the features behind each confirmed drift with TreeSHAP. ASPT applies sequential Monte Carlo stopping to the permutation test, so that each alarm receives a permutation p-value after a small fraction of the maximum number of permutations. EADD was compared with seven unsupervised detectors on synthetic and real-world data streams. It was the only detector to detect all four synthetic drift types (abrupt, gradual, incremental and recurring) without false alarms, and on the real-world streams it detected drift on three of four annotated streams, sooner than D3. An ablation in which D3 was given a LightGBM classifier showed that EADD's advantage on gradual and incremental drift comes from its long-memory reference window rather than from the classifier, that D3's false alarms on stable data were caused by its small windows, and that the permutation test changes the decisions only when the AUC threshold is lowered.
}
